% ============================================================
%  Chapter 6 — Discussion
% ============================================================
\chapter{Discussion}
\label{chap:discussion}

\section{The Central Finding: Weather Is Regime-Dependent}

The thesis's primary empirical contribution is a finding with two parts that
must be understood together. In the stable 2024--25 test period, weather
features add no statistically significant forecast improvement over a
fundamental-only Random Forest ($p = 0.572$, DM = $-0.565$). In the 2022
energy crisis, the same comparison yields $p < 0.001$ (DM = $-13.27$):
weather is highly significant. The role of weather in French day-ahead
electricity price forecasting is therefore not a binary yes/no question but a
function of market regime.

This result reconciles apparently contradictory intuitions. Weather obviously
drives electricity demand and renewable supply, so one would expect it to
matter for price forecasting. Yet in a market where price-setting fundamentals
(TTF gas, nuclear availability) are explicitly modelled, weather information
may be fully subsumed. The resolution is that the degree of subsumption depends
on how tight the coupling between weather and fundamentals is — and this
coupling changes under stress.

\subsection{The Information Redundancy Hypothesis}

We argue that in stable market conditions, weather information is already
embedded in other features of the model through three distinct channels:

\textbf{Load forecast as a temperature proxy.} RTE publishes an hourly load
forecast for D+1 that is computed using a proprietary model incorporating
temperature forecasts \citep{rte2023}. When we include this load forecast in
the feature matrix, we implicitly condition on a temperature-adjusted demand
estimate. The raw temperature variable then adds no marginal information beyond
what is already captured by RTE's own weather-informed demand forecast. The
high correlation between HDD and load forecast ($r > 0.60$, visible in
Figure~\ref{fig:corr_matrix}) confirms this redundancy.

\textbf{TTF gas price as an energy weather proxy.} European gas prices respond
to heating demand aggregated across all gas-consuming countries. A cold spell
in France simultaneously increases French electricity demand and increases gas
demand across Europe, pushing TTF higher. A model that includes TTF therefore
already incorporates a pan-European temperature signal through the lens of the
commodity market. This indirect transmission channel reflects a well-established principle: commodity
markets aggregate weather information efficiently when the weather-commodity
relationship is stable.

\textbf{Nuclear availability dominance.} In France, the availability of the
nuclear fleet explains a disproportionate share of price variance relative to
other European markets. In 2022, nuclear availability fell to $\approx$40\%,
creating price variance that no weather feature can explain (the outages were
driven by corrosion-related technical failures, not weather). With nuclear
availability as an explicit feature, the model's residual variance is already
low in stable conditions, leaving little room for weather variables to
contribute further.

\subsection{Why the Crisis Breaks Redundancy}

The 2022 result ($p < 0.001$) reveals that the information redundancy is
not unconditional. Three mechanisms can explain why weather becomes significant
in a crisis regime:

\textbf{Decoupling of TTF and local temperature.} During the 2021--2023 gas
supply crisis, TTF prices were predominantly driven by geopolitical supply
shocks (Russian gas curtailments) rather than weather. The normal correlation
between cold temperatures and TTF rises was disrupted: prices were high
regardless of the weather. In this environment, the residual temperature effect
on electricity demand could not be subsumed by TTF — it had to be captured
by the raw temperature variable directly. The information redundancy broke down
because the mediating variable (TTF) lost its weather signal.

\textbf{Nuclear constraint amplification.} With nuclear availability at 40\%,
France had limited buffer capacity to absorb demand surges driven by cold
weather. In normal conditions (nuclear at 70\%), a cold spell adds demand that
nuclear can partially accommodate; with nuclear severely constrained, the same
cold spell forces much higher dispatch of expensive gas and oil peakers,
amplifying the price response. This non-linear amplification makes the direct
temperature-price channel visible in the data.

\textbf{Higher signal-to-noise ratio.} The 2022 crisis produced extreme price
levels that created larger absolute prediction errors, but also larger
differences between Model B and Model C: even a small absolute MAE improvement
(0.73 EUR/MWh) is detectable by the DM test at the high absolute price levels
of 2022. The test's statistical power to detect a given proportional effect
is higher when the loss differential has a larger mean.

\subsection{Implications for Forecasting Practice}

This finding has concrete practical implications:

\textit{Stable markets:} Building a rich fundamental model that includes
cross-border flows, fuel prices, and nuclear availability is sufficient to
capture the weather signal implicitly in the French market. The additional cost
and complexity of integrating real-time weather forecasts from ERA5 or NWP
models may not be justified, provided that the load forecast and fuel price
data are available. This lowers the barrier to entry for model builders
without proprietary meteorological data infrastructure.

\textit{Crisis markets:} When TTF-temperature co-movement decouples (e.g.,
during geopolitical supply shocks) or when nuclear availability is severely
constrained, weather features provide statistically and economically significant
value. Market participants entering production forecasting should build
weather-aware models as a contingency, even if weather features are switched off
during normal regimes. A regime-switching model that activates or de-weights
weather features based on nuclear availability and TTF volatility is a
promising direction for future work.

\textit{Cross-market applicability:} This result may not generalise to markets
with higher renewable penetration. In Germany, renewables accounted for over 60\% of electricity generation in 2024
(Fraunhofer ISE, 2025), with wind and solar jointly representing close to half of
all generation, and the direct impact of weather on supply is much larger.
The indirect TTF channel is less likely to fully capture the wind-solar
generation effect in Germany than in France, suggesting that weather features
would be more valuable in the German market context. Testing this hypothesis
in a comparative France-Germany study is a natural extension.

\section{XGBoost Underperformance and Its Implications}

The underperformance of XGBoost (Model D) relative to Random Forest (Models B
and C) is a secondary finding that deserves careful interpretation. With our
hyperparameters, XGBoost achieves MAE of 19.14 EUR/MWh and R$^2$ of 0.659 in
the stable period, compared to approximately 16.94 EUR/MWh and 0.776 for the
Random Forests. In the 2022 crisis, the gap widens further (76.32 vs 66.55
EUR/MWh). Two factors explain this pattern.

First, XGBoost's sequential boosting mechanism assigns higher weight to
poorly-fitted observations in each iteration. In the presence of genuine
outliers driven by extraordinary events (the 2022 nuclear-gas crisis), this
can lead to overfitting on crisis-era residuals at the expense of
normal-market performance. Random Forests average over bootstrap samples,
which is more robust to distributional shifts across regimes.

Second, the hyperparameters used here (500 estimators, $\eta = 0.05$, max
depth 6) were chosen to be reasonable defaults but were not tuned to the
specific loss landscape of our data. The Random Forest's hyperparameters
(500 trees, max depth 10, min leaf 5) were similarly not optimised, but RF
is known to be less sensitive to hyperparameter choice than XGBoost, whose
performance depends more critically on the learning rate, regularisation
terms ($\lambda$, $\alpha$), and subsampling rates. A Bayesian hyperparameter
search with time-series cross-validation could plausibly close or reverse the
performance gap between XGBoost and RF, and is a priority for future work.

The implication for practitioners is that the comparison between RF and
XGBoost reported here should \textit{not} be interpreted as evidence that RF
is inherently superior to boosted trees for electricity price forecasting. The
result reflects our specific configuration. The comparison between Model B and
Model C (RF with vs without weather), however, is robust: both models use the
same architecture and differ only in the feature set, making the DM test a
clean ablation study.

\section{Limitations}
\label{sec:limitations}

\subsection{Static Model and Distribution Shift}

Both RF and XGBoost are trained once on data from 2018--2024 and applied
without retraining to the test period. The training set includes the 2022 crisis
regime, while the test set covers a normalised post-crisis regime. In principle,
the model trained on crisis-level prices (occasionally exceeding 1{,}000
EUR/MWh) may not generalise optimally to a test regime where prices rarely
exceed 200 EUR/MWh. A walk-forward expanding-window retraining scheme — in
which the model is retrained monthly using all data available up to that point
— would provide a more realistic assessment of production performance and is
a prerequisite for operational deployment.

\subsection{Price Spike Forecasting}

Both models systematically underestimate extreme price events. During the
test period, hours with prices above 200 EUR/MWh contribute disproportionately
to RMSE despite representing fewer than 1\% of observations. Point forecasting
models trained on MAE or RMSE loss functions target the conditional mean, which
is pulled toward the bulk of the distribution. Dedicated spike detection
approaches — including two-stage models (a classifier for spike probability
followed by a spike magnitude regression) or quantile regression forests for
prediction intervals — could substantially improve tail performance. This is
particularly important for risk management applications where exposure to
extreme hours drives the economic value of the forecast.

\subsection{Negative Price Dynamics}

Negative prices represent approximately 3--8\% of hours in recent years
(Figure~\ref{fig:price_dist}). Our sMAPE metric handles negative prices
correctly, but the RF models may not accurately represent the specific dynamics
of price formation below zero, where the mechanism shifts from demand-driven
to supply-push (inflexible generators paying consumers to absorb excess
generation). A dedicated modelling approach for the negative-price regime is
a natural extension.

\subsection{Missing EU ETS Carbon Allowances}

EU ETS carbon allowance prices (EUA) were not included as a standalone feature.
EUA prices directly raise the operating cost of gas and coal generation through
the carbon component of the generation cost. Their exclusion may reduce the
model's accuracy during periods of high carbon price volatility, which have
coincided with some of the most extreme electricity price episodes (2021--2022).
The TTF-coal spread partially proxies for EUA (as all three series moved
together during the crisis), but a more complete fundamental model would include
EUA explicitly.

\subsection{Single Out-of-Sample Regime}

The main test period (May 2024--April 2025) covers only one market regime.
The 2022 robustness exercise provides an important second data point, but does
not constitute a full walk-forward evaluation. The reliability of the
statistical conclusions would be strengthened by testing on additional
out-of-sample windows covering different nuclear availability levels, renewable
penetration levels, and geopolitical conditions.

\section{Future Research Directions}

\begin{enumerate}
  \item \textbf{Regime-switching weather model:} Building an explicit
  regime-switching mechanism that activates weather features when nuclear
  availability falls below a threshold (e.g., 0.55) or when TTF volatility
  exceeds a threshold, combining the efficiency of the fundamental model in
  stable regimes with the accuracy of the weather-augmented model in stress
  regimes.

  \item \textbf{Probabilistic forecasting:} Replace point forecasts with
  prediction intervals or full predictive distributions using Quantile
  Regression Forests or conformal prediction. Probabilistic forecasts are
  increasingly demanded by market participants for option pricing and risk
  management.

  \item \textbf{Comparative France--Germany study:} Test whether weather
  features are more significant in Germany, where renewable penetration is
  substantially higher and the indirect TTF-weather channel is less dominant.
  Such a comparison would directly test our information redundancy hypothesis
  in a different market structure.

  \item \textbf{XGBoost hyperparameter optimisation:} A systematic Bayesian
  search with temporal cross-validation, using a custom time-aware split
  strategy (e.g., 6-fold expanding window), could close the performance gap
  between XGBoost and RF.

  \item \textbf{Transformer and LSTM architectures:} Long Short-Term Memory
  networks and Transformer-based models have demonstrated promise on energy
  time series with long-range dependencies \citep{lago2018}. Comparing these
  to our RF baseline within the same experimental framework — including the
  same DM test procedure — would be a natural extension.
\end{enumerate}
